\section{文档概述}

\subsection{编写目的}

本文建立「寻迹」校园失物招领平台的需求基线，说明项目要解决的问题、建设目标、系统范围、用户角色、功能需求、性能需求、可维护性需求、其他质量要求和业务规则。报告用于统一团队对「系统应当提供什么能力」的理解，并作为系统建模、架构设计、开发和测试工作的共同输入。

本文只定义需求及其验收条件，不评价代码实现状态，也不记录测试执行结果。文中「应」表示必须满足的系统行为，「不得」表示禁止出现的行为；P0、P1、P2 分别表示核心需求、重要增强需求和扩展需求。

\subsection{需求依据}

\begin{table}[H]
  \centering
  \caption{需求来源与用途}
  \begin{tabularx}{\textwidth}{L{3.0cm} Y Y}
    \toprule
    来源 & 主要内容 & 在本报告中的用途 \\
    \midrule
    课程期末要求 & 项目目标、功能、性能、可维护性等报告要求 & 确定报告范围和质量要求 \\
    校园失物招领场景 & 信息分散、检索困难、冒领和隐私保护问题 & 提炼业务目标、角色和使用场景 \\
    团队需求基线 & P0/P1/P2 范围、业务规则和验收口径 & 形成需求编号、优先级与边界 \\
    枚举与交互规范 & 角色、状态和跨模块业务语义 & 保持术语、状态和规则一致 \\
    \bottomrule
  \end{tabularx}
\end{table}

\subsection{核心术语}

\begin{table}[H]
  \centering
  \caption{核心业务术语}
  \begin{tabularx}{\textwidth}{L{2.8cm} Y}
    \toprule
    术语 & 定义 \\
    \midrule
    失物信息 & 失主发布的寻物记录，描述丢失物品、时间、地点和特征 \\
    招领信息 & 拾获者或后勤人员发布的拾获记录，描述物品特征、保管方式和验证要求 \\
    匹配结果 & 系统根据文本、类别、地点、时间等信息形成的候选关联及评分 \\
    认领申请 & 用户针对招领信息发起的身份与物品归属验证流程 \\
    交接 & 认领审核通过后，失主与拾获者确认物品归还的业务过程 \\
    敏感物品 & 证件或包含个人信息、需要限制图片和详情访问的物品 \\
    \bottomrule
  \end{tabularx}
\end{table}

\section{项目背景与目标}

\subsection{问题陈述}

校园失物信息通常分散在聊天群、公告栏和个人社交账号中，信息格式不统一、检索效率低、有效期难管理。传统的「看到信息后直接联系」方式还缺少统一的归属验证和交接记录，容易产生冒领、隐私泄露和责任边界不清等问题。

「寻迹」面向校内失物招领场景，将失物与招领信息转化为结构化记录，以用户端承载发布、搜索、匹配、认领和交接，以管理端承载认证、审核和治理，并通过智能服务辅助分类、敏感检测和候选匹配。

\subsection{项目目标}

\begin{table}[H]
  \centering
  \caption{项目目标}
  \begin{tabularx}{\textwidth}{L{1.6cm} L{3.2cm} Y}
    \toprule
    编号 & 目标 & 目标说明 \\
    \midrule
    G-01 & 统一信息管理 & 建立结构化的失物与招领信息库，使信息能够发布、查询、筛选和持续维护 \\
    G-02 & 提升找回效率 & 根据物品特征、地点和时间生成可解释的候选匹配，并及时通知相关用户 \\
    G-03 & 保障认领安全 & 根据物品敏感程度和用户信誉设置分级验证，通过审核与双确认交接减少冒领 \\
    G-04 & 保护个人隐私 & 限制证件、认领凭证和敏感图片的访问范围，避免公开暴露个人信息 \\
    G-05 & 支持校园治理 & 为后勤和管理员提供认证、内容审核、举报处置、公告、统计和操作追踪能力 \\
    \bottomrule
  \end{tabularx}
\end{table}

\section{干系人与使用需求}

\begin{longtable}{L{2.4cm} L{3.2cm} L{5.0cm} L{3.6cm}}
  \caption{角色、使用目标与权限边界}\label{tab:roles}\\
  \toprule
  角色 & 使用目标 & 主要需求 & 权限边界 \\
  \midrule
  \endfirsthead
  \toprule
  角色 & 使用目标 & 主要需求 & 权限边界 \\
  \midrule
  \endhead
  访客 & 建立合法会话 & 注册、登录、找回访问凭据 & 不得访问用户业务数据和管理功能 \\
  普通用户（USER） & 找回失物或归还拾获物 & 认证、发布、搜索、匹配、认领、交接、通知、积分和举报 & 只能管理本人发布及本人参与的业务记录 \\
  失主 & 找到并安全取回物品 & 发布失物、订阅匹配、查看候选、发起认领和确认收到 & 不得认领本人发布的招领信息 \\
  拾获者 & 发布并安全归还物品 & 发布招领、设置验证问题、审核认领、安排和确认交接 & 只能审核本人所发布招领的认领申请 \\
  后勤/安保（STAFF） & 登记和代管拾获物品 & 批量登记、保管说明、认领审核和交接确认 & 仅处理职责范围内的代管记录 \\
  管理员（ADMIN） & 审核、治理和运营 & 认证审批、内容审核、举报处置、公告、统计与日志查询 & 管理操作必须经过授权并保留操作记录 \\
  AI 辅助服务 & 提供辅助判断 & 分类、敏感检测、文本相似度和匹配评分 & 只返回计算结果，不决定最终认领归属 \\
  \bottomrule
\end{longtable}

\section{系统范围}

\subsection{范围内能力}

系统范围包括账户与校园认证、失物和招领发布、图片管理、搜索筛选、候选匹配、分级认领、申诉、交接确认、站内通知、信誉积分、举报治理、公告、统计和操作记录。普通用户通过移动端响应式网页使用业务功能，管理员通过 PC 管理端进行审核与治理。

\subsection{系统边界}

系统与用户、管理员、对象存储和 AI 辅助服务发生交互。业务后端负责身份校验、业务规则和状态一致性；对象存储负责保存图片与凭证；AI 辅助服务只提供分类、敏感判断和相似度计算，服务不可用时不得阻断基础发布，也不得降低隐私保护等级。

\subsection{范围外能力}

下列能力不属于本需求基线：与学校统一认证、教务系统或真实身份库直接对接；真实短信、支付和物流服务；实时聊天与完整消息中台；多租户和分校区独立运营；在线模型训练及商业级图像检索平台。范围外能力不得成为 P0 主流程的前置条件。

\section{功能需求}

\subsection{功能需求总览}

\small
\begin{longtable}{L{1.4cm} L{3.0cm} L{6.8cm} L{2.4cm}}
  \caption{功能需求基线}\label{tab:functional}\\
  \toprule
  编号 & 名称 & 需求说明 & 优先级 \\
  \midrule
  \endfirsthead
  \toprule
  编号 & 名称 & 需求说明 & 优先级 \\
  \midrule
  \endhead
  FR-01 & 账户、资料与认证 & 系统应支持普通用户注册登录、资料维护、校园认证申请与审批，并为管理员提供独立登录入口 & P0 \\
  FR-02 & 失物发布与检索 & 用户应能发布、查询、编辑、关闭失物信息，并管理本人历史记录 & P0 \\
  FR-03 & 招领发布与保护 & 拾获者应能发布和维护招领信息、设置验证问题；系统应限制敏感物品信息访问 & P0/P1 \\
  FR-04 & 搜索与候选匹配 & 系统应支持多条件搜索，并根据可用信息生成、解释和展示候选匹配 & P0/P1 \\
  FR-05 & 认领、申诉与交接 & 系统应支持分级验证、认领审核、申诉、交接安排和双方确认 & P0/P2 \\
  FR-06 & 信誉积分 & 系统应记录用户积分及变动原因，并依据业务事件进行一致、不可重复的调整 & P1 \\
  FR-07 & 通知与公告 & 系统应为匹配、认领、审核、交接和积分事件生成站内通知，并展示有效公告 & P0 \\
  FR-08 & 后台审核与治理 & 管理员应能处理认证、内容、举报、申诉、用户状态、公告、统计和操作记录 & P0/P1 \\
  \bottomrule
\end{longtable}
\normalsize

\subsection{FR-01：账户、资料与校园认证}

\begin{enumerate}
  \item 系统应支持使用手机号、验证码、密码和昵称注册普通用户，手机号在系统内必须唯一。
  \item 普通用户应能使用手机号与密码或验证码登录；管理员应使用独立账号和密码登录管理端。
  \item 登录凭据必须设置有效期；账号被禁用或注销后，原有登录凭据不得继续访问受保护资源。
  \item 用户应能查看和修改允许变更的个人资料，并查看校园认证状态。
  \item 用户应能提交校园编号、真实姓名和证件图片发起认证；同一用户同时最多存在一项待处理申请。
  \item 管理员应能批准或驳回待处理认证，处理结果应通知申请人并记录操作人、时间和结果。
\end{enumerate}

\subsection{FR-02：失物发布与生命周期}

\begin{enumerate}
  \item 失物信息必须包含名称、类别、丢失时间区间和地点；描述和图片可选，图片数量不得超过 5 张。
  \item 发布者应能选择是否接收匹配通知，并在提交前预览所填写的信息。
  \item 用户应能按状态查看本人发布的失物，修改仍处于可编辑状态的记录，并主动关闭不再查找的记录。
  \item 失物进入已找回或已关闭状态后，不应继续参与新的活动匹配。
  \item 记录修改后，系统应依据新信息重新计算候选匹配，旧的无效候选不得继续用于发起认领。
\end{enumerate}

\subsection{FR-03：招领发布与敏感保护}

\begin{enumerate}
  \item 招领信息必须包含名称、类别、拾获时间、拾获地点、保管方式和联系偏好；描述和图片可选，图片数量不得超过 5 张。
  \item 发布者可设置最多 3 个验证问题。面向认领者的页面只能展示问题，不得展示答案关键词或评分规则。
  \item 系统应根据物品类别和敏感检测结果标记敏感物品。证件和敏感原图只能向发布者、合法认领参与者或管理员展示。
  \item 拾获者应能查看、编辑和关闭本人发布的招领；存在活动认领时，关闭操作必须同步终止相关流程并通知参与者。
  \item STAFF 角色应支持一次登记不超过 50 条招领信息，并对每条失败记录给出可定位的校验结果。
\end{enumerate}

\subsection{FR-04：搜索与候选匹配}

\begin{enumerate}
  \item 用户应能按关键词、类别、地点、事件时间范围和业务状态搜索失物与招领，并按创建时间或事件时间排序。
  \item 搜索结果应分页展示，默认排除业务终态和无权访问的内容；空结果、加载失败和无访问权限必须使用不同提示。
  \item 发布或编辑物品后，系统应在不阻塞发布响应的前提下生成候选匹配。
  \item 匹配应综合文本、地点、时间及可用的图像特征，并向用户展示总分和主要匹配依据。
  \item 匹配分数按可用维度归一计算：
  \[
    S=\frac{\sum_{i\in A}w_i s_i}{\sum_{i\in A}w_i},\qquad
    (w_{image},w_{text},w_{location},w_{time})=(0.4,0.3,0.2,0.1).
  \]
  当某一维度不可用时，该维度不得以默认高分参与计算；总分达到 70 方可形成活动匹配。
  \item AI 辅助服务不可用时，系统应使用文本、地点和时间规则生成基础候选；降级不得放宽敏感信息访问限制。
\end{enumerate}

\subsection{FR-05：认领、申诉与交接}

\begin{enumerate}
  \item 用户应能从招领详情或合法匹配上下文发起认领；系统必须阻止用户认领本人发布的招领。
  \item 验证等级由物品类别和认领者信誉共同确定：普通物品默认一级问答验证，电子设备默认二级问答加凭证验证，证件固定三级凭证与人工审核；低信誉用户应提高验证等级或禁止发起认领。
  \item 问答必须完整提交，平均关键词命中率达到 0.6 方可通过；失败结果不得泄露逐题答案或分数。
  \item 同一认领者针对同一招领连续失败后应进入 5 分钟冷却期，滚动 24 小时内最多失败 3 次。
  \item 同一招领同时最多存在一个活动认领。并发申请时，系统必须保证只有一个申请能够占用该招领。
  \item 拾获者应能审核认领申请；认领被拒绝后，认领者应能提交申诉，由管理员作出处理决定。
  \item 审核通过后，双方应能约定当面交接或指定代收点。失主确认收到、拾获者确认归还后，系统应同步更新认领、招领和关联失物状态。
  \item 重复创建同一交接或重复确认已处理动作时，系统不得生成重复记录或重复积分。
\end{enumerate}

\subsection{FR-06：信誉积分}

系统应向用户展示当前积分和积分流水。校园认证、成功交接、有效举报和违规处置等事件可触发积分变化，每条流水必须说明时间、原因和关联业务。同一业务事件不得重复计分，积分规则变更后不得篡改既有流水。

\subsection{FR-07：通知与公告}

系统应为匹配生成、认领申请、审核结果、交接提醒、积分变化和系统公告生成站内通知。用户应能查看通知列表与未读数量、标记已读，并从通知进入对应业务页面。通知必须指向用户有权访问的对象；关联对象不可用时，应说明原因而不是进入错误页面。

\subsection{FR-08：后台审核与治理}

管理员应能审批校园认证和内容、处理认领申诉与用户举报、调整用户状态、发布和下线公告、查看统计看板及查询操作记录。举报处置应关联举报对象、处理意见和处理人；涉及下架、终止认领或积分变化时，应通知受影响用户。普通用户和 STAFF 不得访问管理员接口。

\section{性能需求}

\subsection{度量口径}

性能指标以 50 名并发活跃用户、10 万条物品记录、每页 20 条结果为基准负载。接口时间从服务端收到请求计至返回完整响应，页面时间从导航开始计至关键内容可见，匹配时间从任务被接受计至候选结果可查询。除特别说明外，响应时间采用 P95 统计口径。

\small
\begin{longtable}{L{1.5cm} L{2.6cm} L{5.2cm} L{4.8cm}}
  \caption{性能需求}\label{tab:performance}\\
  \toprule
  编号 & 场景 & 适用条件 & 指标要求 \\
  \midrule
  \endfirsthead
  \toprule
  编号 & 场景 & 适用条件 & 指标要求 \\
  \midrule
  \endhead
  PERF-01 & 普通业务接口 & 不含文件上传和 AI 计算的查询、创建与状态操作 & P95 响应时间不超过 500 ms，错误率不高于 1\% \\
  PERF-02 & 物品搜索 & 10 万条物品记录、组合两个筛选条件并按时间排序 & P95 响应时间不超过 1 s \\
  PERF-03 & 关键页面首屏 & 20 Mbps 网络、主流移动设备、静态资源已正常部署 & 登录后首页、搜索页和详情页关键内容在 1.5 s 内可见 \\
  PERF-04 & 候选匹配 & AI 服务可用、单批不超过 20 个并行匹配任务 & 95\% 的单次匹配在 3 s 内形成可查询结果 \\
  PERF-05 & 状态通知 & 匹配、认领、审核、交接和积分状态写入成功后 & 站内通知在 2 s 内可查询，未读数同步更新 \\
  PERF-06 & 并发认领 & 10 个用户同时认领同一招领信息 & 请求均在 2 s 内得到明确结果，且最多产生一个活动认领 \\
  \bottomrule
\end{longtable}
\normalsize

\section{可维护性需求}

\small
\begin{longtable}{L{1.7cm} L{3.0cm} L{8.2cm}}
  \caption{可维护性需求}\label{tab:maintainability}\\
  \toprule
  编号 & 类别 & 需求说明 \\
  \midrule
  \endfirsthead
  \toprule
  编号 & 类别 & 需求说明 \\
  \midrule
  \endhead
  MAINT-01 & 模块边界 & 用户、物品、匹配、认领、积分、通知和管理模块应职责单一；跨模块调用必须经过公开服务接口，不得绕过业务层直接访问其他模块的数据层 \\
  MAINT-02 & 接口契约 & 请求、响应、错误码和枚举应有统一定义；接口字段变化时，接口文档、前端调用方和相关测试必须同步更新 \\
  MAINT-03 & 编码规范 & Python 和 TypeScript 代码应使用明确类型、统一命名和自动格式检查；业务规则不得散落在路由或页面组件中 \\
  MAINT-04 & 配置管理 & 数据库地址、服务地址、密钥、存储配置和功能开关必须从环境配置读取，不得硬编码或提交真实凭据 \\
  MAINT-05 & 数据演化 & 数据库结构变化必须通过有序迁移脚本管理；每次迁移应说明升级路径、依赖条件和恢复方法 \\
  MAINT-06 & 自动化检查 & 核心业务规则、权限边界和状态流转必须具备自动化测试；合并前应执行格式、静态检查、类型检查和相关测试 \\
  MAINT-07 & 文档同步 & API、枚举、数据库字段或业务流程变化时，需求、架构、接口和测试文档中的对应内容必须在同一变更中更新 \\
  MAINT-08 & 问题定位 & 服务应输出包含时间、级别、请求标识和业务上下文的结构化日志；敏感信息不得写入日志 \\
  \bottomrule
\end{longtable}
\normalsize

\section{其他非功能需求}

\small
\begin{longtable}{L{1.6cm} L{2.4cm} L{8.9cm}}
  \caption{其他非功能需求}\label{tab:other-nfr}\\
  \toprule
  编号 & 类别 & 需求说明 \\
  \midrule
  \endfirsthead
  \toprule
  编号 & 类别 & 需求说明 \\
  \midrule
  \endhead
  NFR-01 & 安全性 & 密码必须采用带盐慢哈希保存；受保护操作必须校验登录状态和角色；令牌、密钥和验证码不得以明文日志输出 \\
  NFR-02 & 可靠性 & 发布、认领和交接等跨实体状态变化必须保持一致；临时任务失败后应自动重试，重复请求不得产生重复副作用 \\
  NFR-03 & 易用性 & 用户端发布、搜索、认领和交接等关键任务应在不超过 3 个主要步骤内到达目标；加载、空数据、失败和无权限状态必须清晰区分 \\
  NFR-04 & 兼容性 & 用户端应适配 360--1440 px 宽度，并支持主流 Chromium、Firefox 和 Safari 最近两个主要版本；管理端应支持常用桌面分辨率 \\
  NFR-05 & 隐私性 & 校园证件、认领凭证和敏感物品原图必须私有保存，只向授权参与者或管理员提供短时访问权限 \\
  NFR-06 & 可追溯性 & 认证审批、内容审核、举报处置、积分变化、关键状态变化和交接确认必须记录操作者、角色、对象、动作与时间 \\
  NFR-07 & 可部署性 & 系统应提供环境配置模板、数据库迁移入口、服务健康检查和可重复的部署流程；缺少必要密钥时不得以生产模式启动 \\
  NFR-08 & 可恢复性 & 业务数据和对象文件应支持定期备份；恢复流程应能够校验数据库版本、文件引用和核心业务记录的一致性 \\
  \bottomrule
\end{longtable}
\normalsize

\section{业务规则与数据约束}

\subsection{状态流转规则}

\begin{table}[H]
  \centering
  \caption{核心实体状态规则}
  \begin{tabularx}{\textwidth}{L{2.4cm} Y Y}
    \toprule
    实体 & 主状态路径 & 关键约束 \\
    \midrule
    失物 & 寻物中\,$\rightarrow$\,已找回/已关闭 & 终态不再产生活动匹配；活动认领存在时不得直接关闭 \\
    招领 & 待认领\,$\rightarrow$\,认领中\,$\rightarrow$\,已归还；也可关闭 & 同时最多一个活动认领；关闭时必须处理关联认领和匹配 \\
    匹配 & 新匹配\,$\rightarrow$\,已读\,$\rightarrow$\,已认领/已失效 & 同一失物和招领组合只有一个活动匹配；上游信息失效时匹配必须过期 \\
    认领 & 待验证\,$\rightarrow$\,已通过/已拒绝/申诉中\,$\rightarrow$\,已交接 & 被上游关闭或治理终止时应进入「已终止」，不得与普通拒绝混用 \\
    公告 & 草稿\,$\rightarrow$\,已发布\,$\rightarrow$\,已下线 & 下线公告不得继续出现在用户公告列表 \\
    \bottomrule
  \end{tabularx}
\end{table}

\subsection{一致性与幂等规则}

认领占用、认领审核、交接确认和积分变化必须保持业务一致性。同一招领不得出现两个活动认领；双方确认交接后，认领、招领及合法关联失物应同步进入终态。网络重试或重复点击不得造成重复交接、重复通知、重复积分或相互矛盾的状态。

\subsection{数据与隐私规则}

JSON 字段采用 camelCase，数据库列采用 snake\_case，业务枚举必须使用统一字典。图片记录应保存稳定资产标识，不得把永久公开地址作为业务数据。认领回答应保留问题文本快照，以便在问题修改后还原当时的验证上下文。核心业务记录原则上采用状态关闭，不得因物理删除产生悬空匹配、通知或操作记录。

\section{需求优先级}

\subsection{P0 核心需求}

\small
\begin{longtable}{L{1.5cm} L{4.0cm} L{4.4cm} L{4.0cm}}
  \caption{P0 核心需求}\label{tab:p0}\\
  \toprule
  编号 & 功能 & 业务价值 & 验收关注点 \\
  \midrule
  \endfirsthead
  \toprule
  编号 & 功能 & 业务价值 & 验收关注点 \\
  \midrule
  \endhead
  P0-01 & 注册、登录与资料 & 建立可识别的用户会话 & 普通用户和管理员入口分离，账号状态有效 \\
  P0-02 & 校园认证申请与审批 & 建立校园身份可信度 & 申请、审批、结果通知和敏感图片权限明确 \\
  P0-03 & 失物发布 & 形成可检索的寻物信息 & 必填字段、图片数量和生命周期规则明确 \\
  P0-04 & 招领发布 & 形成可验证的拾获信息 & 保管方式、验证问题和敏感保护明确 \\
  P0-05 & 搜索与筛选 & 帮助用户主动发现候选 & 类别、地点、时间、状态和排序组合有效 \\
  P0-06 & 匹配结果展示 & 降低人工比对成本 & 候选阈值、解释信息和降级行为明确 \\
  P0-07 & 认领与一级验证 & 建立基础防冒领流程 & 归属校验、完整问答、失败冷却和互斥占用有效 \\
  P0-08 & 双确认交接 & 形成归还闭环 & 双方确认后关联状态一致且重复操作幂等 \\
  P0-09 & 站内通知 & 及时传递关键业务事件 & 未读状态、业务跳转和权限边界正确 \\
  P0-10 & 后台基础审核 & 支持校园内容治理 & 认证、内容、举报和操作记录可管理 \\
  \bottomrule
\end{longtable}
\normalsize

\subsection{P1 与 P2 需求}

P1 包括综合匹配评分、敏感物品受控展示、信誉积分、举报处置、数据统计和 STAFF 批量录入，用于增强平台完整性。P2 包括二级/三级凭证审核、认领申诉、外部通知、长期未领清单和更强的图像或 OCR 能力，用于扩展复杂场景。P1、P2 不得改变 P0 主流程的状态语义和隐私边界。

新增需求应先分析对数据库、枚举、接口、页面、模型和验收条件的影响；影响核心流程的变更必须更新需求基线后再进入设计。

\section{验收条件}

\subsection{核心业务场景}

\begin{enumerate}
  \item 普通用户应能够注册或使用密码登录，并查看个人资料和认证状态。
  \item 用户 A 应能够发布一条订阅匹配的失物，用户 B 应能够发布一条带验证问题的招领。
  \item 系统应在发布响应之后生成可解释候选，并向订阅用户发送站内通知。
  \item 用户 A 应能够从合法匹配上下文发起认领，系统应根据物品类别和信誉执行正确验证等级。
  \item 拾获者应能够审核认领，审核通过后双方应能够创建交接并分别确认。
  \item 双方确认后，认领、招领和关联失物状态应保持一致，积分、通知和操作记录不得重复。
  \item 管理员应能够处理认证、内容和举报，并查看统计与关键操作记录。
\end{enumerate}

\subsection{异常与安全场景}

系统必须能够区分并正确处理未登录访问、越权访问、重复提交、并发认领、问答缺失、连续失败、敏感图片无权访问、AI 服务超时、任务中断、上游记录关闭、并发审批和并发双确认。错误页面不得将网络异常显示为「暂无数据」，任何辅助服务故障不得解除敏感保护或自动通过证件认领。

\subsection{质量验收口径}

性能验收应按照第 6 章的负载、数据规模、计时边界和 P95 口径执行。可维护性验收应检查模块依赖、接口契约、配置外置、迁移管理、自动化检查和文档同步机制。安全与隐私验收应覆盖凭据保护、角色权限、敏感图片访问和关键操作追踪。

\section{需求内部追踪}

\begin{table}[H]
  \centering
  \caption{项目目标与需求追踪}
  \begin{tabularx}{\textwidth}{L{1.6cm} L{3.0cm} L{4.2cm} Y}
    \toprule
    目标 & 功能需求 & 质量需求 & 对应验收场景 \\
    \midrule
    G-01 & FR-02、FR-03、FR-08 & PERF-02、MAINT-02 & 发布、检索、维护和治理结构化信息 \\
    G-02 & FR-04、FR-07 & PERF-02、PERF-04、PERF-05 & 生成可解释候选并通知相关用户 \\
    G-03 & FR-01、FR-05、FR-06 & NFR-01、NFR-02、NFR-06 & 分级验证、互斥认领和双确认交接 \\
    G-04 & FR-01、FR-03、FR-05 & NFR-01、NFR-05 & 未授权用户无法访问证件和敏感原图 \\
    G-05 & FR-08 & MAINT-08、NFR-06、NFR-07 & 管理员完成审核、处置、统计和追踪 \\
    \bottomrule
  \end{tabularx}
\end{table}

\section{结论}

本报告从校园失物招领问题出发，明确了「寻迹」的五项目标、七类干系人、系统边界、八组功能需求以及性能、可维护性、安全、可靠性、易用性、隐私和部署要求。业务规则、需求优先级、验收条件与内部追踪共同构成统一需求基线，为后续系统建模、架构设计、开发和测试提供明确输入。
